\section{Discussion}
\label{sec:discussion}

Taken together, the empirical results support a layered interpretation of Brazilian equity dependencies. The first layer is the Market Mode. It is visible in the moderate positive average correlation ($\bar{\rho} \approx 0.24$), the increase in rolling average correlation during stress periods and the very large first eigenvalue ($\lambda_1 = 21.65$, approximately $37.3\%$ of the spectrum). This mode is economically meaningful because it captures broad market synchronization driven by domestic monetary policy, exchange-rate dynamics and global commodity cycles. It is also methodologically important because it can dominate raw networks and obscure weaker local structures.

The second layer is group and sector structure. Within-sector correlations are statistically significantly higher than between-sector correlations ($p \approx 1.33 \times 10^{-24}$), and the RMT decomposition identifies four additional eigenvalues above the random upper edge. Eigenvector loadings and group-mode matrices suggest that commodity, financial, telecommunications and utility-related structures contribute to the non-random part of the spectrum. These structures are not as strong as the market mode, but they are precisely the level where network filtering becomes informative.

The third layer is noisy residual dependence. Not every correlation is meaningful, and not every edge retained by a graph should be interpreted economically. This is why the paper combines RMT, heatmaps, dendrograms and network metrics. A result is more credible when it appears consistently across several representations rather than only in one visualization.

The network results demonstrate that market-mode removal changes topology in a structured way. Original networks have stronger average edge correlations, while group-mode networks have much lower average edge correlations. However, the group-mode PMFG has higher clustering (0.6653 versus 0.5273). This combination is a central finding: lower average edge strength does not mean absence of structure. It means that the dominant common component has been removed and localized mesoscopic organization becomes easier to observe. The PMFG is particularly suited for this analysis because it preserves richer topology than the MST, retaining 168 edges with 166 triangles and 55 4-cliques that allow hub-rank shifts, clustering changes and subsector dependency maps to be studied in detail.

Clustering summarizes hierarchical organization, but it does not identify graph-theoretic hubs, bridges or local cycles. The transition from dendrograms to filtered financial networks reveals complementary information. The hub reorganization after RMT filtering is particularly striking: original hubs such as GGBR4 and BBDC4 reflect market-wide correlation strength and broad exposure, while group-mode hubs such as GUAR3 represent localized bridges after the common component is removed. This suggests that centrality in financial networks has different economic interpretations depending on whether the market mode is included or excluded.

The forecasting results should be interpreted with appropriate caution. They do not show that machine learning uniformly dominates classical volatility models. Instead, they show that network-augmented models can approach strong baselines and that PMFG and MST variables have nonzero importance in Random Forest models. The persistence of GARCH(1,1), with $\alpha + \beta$ consistently above 0.98, confirms that autoregressive volatility dynamics remain the dominant forecasting signal for Brazilian equities. Network features complement rather than replace this signal. The appropriate interpretation is incremental: topology adds structural information from the cross-sectional geometry of the market to the temporal dynamics of individual asset variance.

The broader contribution is methodological. The paper connects econophysics filtering, financial networks and volatility forecasting in a single reproducible pipeline for B3 equities. By documenting each step transparently and evaluating network features against demanding econometric benchmarks, the analysis provides a template that can be extended to other emerging markets, different time scales and more sophisticated forecasting architectures.
